\chapter{Modeling and the Question Phase}

\otbookepigraph{A relationship becomes secure only after it acquires an internal reason.}{PLDS Principle}

\section{The Structure of Rotation}

In Spong Chapter 2, we encounter the \otindex{Rotation Matrix}. Instead of seeing it as a collection of sines and cosines, we \otindex{Model} it as a transformation that preserves the dot product.

\begin{intuitionbox}[Necessity of Orthogonality]
    A rotation matrix \(R\) must satisfy \(R^T R = I\) because the lengths of vectors and the angles between them must not change after rotation.
\end{intuitionbox}

\section{Symbolic Clarity}

Dense texts often reuse symbols. By using the \acr{plds} key-based approach, we distinguish between a \sym{joint-angle} (\(\theta\)) and the \sym{temperature} (\(\theta\)) of the actuator.

\begin{examplebox}[DH Matrix Construction]
    The homogeneous transformation matrix is a composition of:
    \begin{enumerate}
        \item Rotation around \(z\) by \sym{joint-angle}.
        \item Translation along \(z\).
        \item Translation along \(x\).
        \item Rotation around \(x\).
    \end{enumerate}
\end{examplebox}

\section{Attacking the Arbitrary}

When reading Spong's derivation of the \otindex{Euler-Lagrange} equations, don't just follow the algebra. Use the \textit{Strong Question} technique:
\begin{itemize}
    \item \textbf{Weak:} Why is the kinetic energy subtracted from potential?
    \item \textbf{Stronger:} How does the \otindex{Lagrangian} formulation represent the constraints of the system without requiring explicit constraint forces?
\end{itemize}
